\section{A finite tensor obstruction}\label{sec:tensor}

This section proves an algebraic obstruction to the formal data that will
arise from the two vertex pieces.  The formal parametrizations may involve
every edge block; each potential depends only on its own block.
All parameter rings below retain every mixed square-free parameter monomial.

For a characteristic-zero field \(K\), write \(K\langle h\rangle\) for the
ring of algebraic power series at the origin, identified with the
henselization of \(K[h]_{(h)}\) inside \(K[[h]]\).  An algebraic germ can
therefore be represented by a regular function on a pointed étale
neighborhood.  Gradients in this section always differentiate internal
coordinates, keeping the external parameters fixed.

\begin{theorem}[Finite tensor obstruction]\label{thm:tensor-obstruction}
Let \(h=(h_1,\ldots,h_n)\) be a finite list of coordinate blocks over a
characteristic-zero field \(K\), of total dimension \(N\).  Let
\(I\subset\{1,\ldots,n\}\) have cardinality \(m\geq5\), and let
\[
 F_i\in K\langle h_i\rangle\quad(1\leq i\leq n),
 \qquad b_i\in K\langle h_i\rangle\quad(i\in I).
\]
Put
\[
 T=K[t_i:i\in I]/(t_i^2:i\in I),\qquad
 S_0(h)=\sum_{i=1}^nF_i(h_i),\qquad
 S_t(h)=S_0(h)+\sum_{i\in I}t_i b_i(h_i).
\]
Suppose there are formal parametrizations
\[
 H_L(x,t)\in (x,t)\,T[[x]]^N,\qquad
 H_R(y)\in (y)\,K[[y]]^N,
\]
where \(x\) has \(a\) coordinates and \(y\) has \(N-a\) coordinates, with
the following properties.
\begin{enumerate}[label=\textup{(\roman*)}]
 \item The central tangent maps
 \[
  \partial_xH_L(0,0):K^a\longrightarrow K^N,\qquad
  \partial_yH_R(0):K^{N-a}\longrightarrow K^N
 \]
 are injective and have complementary images.
 \item The identities
 \[
     S_t(H_L(x,t))=0,\qquad S_0(H_R(y))=0
 \]
 hold in \(T[[x]]\) and \(K[[y]]\), respectively.
 \item For every three distinct \(i,j,k\in I\),
 \begin{align}
 t_i t_j t_k
 &\in\bigl((\partial_{h_\nu}S_t)(H_L(x,t)):
                    1\leq\nu\leq N\bigr)\subset T[[x]],
       \label{eq:tensor-source-membership}\\
 (b_i b_j b_k)(H_R(y))
 &\in\bigl((\partial_{h_\nu}S_0)(H_R(y)):
                    1\leq\nu\leq N\bigr)\subset K[[y]].
       \label{eq:tensor-target-membership}
 \end{align}
\end{enumerate}
Then these data do not exist.
\end{theorem}

Condition (i) makes the parametrizations smooth formal graph embeddings,
including over the full ring \(T\).  Nilpotent constant displacements of
\(H_L\) are allowed.  The ideals in
\eqref{eq:tensor-source-membership}--\eqref{eq:tensor-target-membership}
are the actual gradient ideals; no passage to their radicals is made.
There is no flatness assumption on their quotients.

The later construction supplies the hypotheses as follows; all triples
in the table have distinct active indices.
\begin{center}
\begin{tabular}{@{}
 >{\raggedright\arraybackslash}p{0.25\linewidth}
 >{\raggedright\arraybackslash}p{0.30\linewidth}
 >{\raggedright\arraybackslash}p{\dimexpr0.45\linewidth-4\tabcolsep\relax}
 @{}}
\toprule
Input & Source of the data & Later verification \\
\midrule
Separate algebraic \(F_i(h_i)\) and \(b_i(h_i)\)
 & Individual marked edge models and polynomial perturbations
 & Theorem~\ref{thm:edge-algebraicity} and
   Section~\ref{subsec:vertex-polynomial} \\[4pt]
Complementary central tangents
 & Restriction maps after fixing the prescribed evaluations
 & Lemma~\ref{lem:vertex-linear} \\[4pt]
Exact vanishing \(S_t|_L=0\) and \(S_0|_R=0\)
 & Stokes and section corrections preserving the residual quotient
 & Lemmas~\ref{lem:vertex-offshell-identities},
   \ref{lem:vertex-correction}, and
   \ref{lem:tensor-nilpotent-correction} \\[4pt]
Every source triple \(t_i t_j t_k\) in the restricted gradient ideal on \(L\)
 & Boundary laws at \(A\)
 & Lemma~\ref{lem:vertex-triple-ideals} \\[4pt]
Every target triple \(b_i b_j b_k\) in the restricted gradient ideal on \(R\)
 & Boundary laws at \(B\) and polynomial holonomy weights
 & Lemma~\ref{lem:vertex-triple-ideals} and
   Section~\ref{subsec:vertex-polynomial} \\
\bottomrule
\end{tabular}
\end{center}

The threshold \(m\geq5\) comes from two counts.  The target triple
identities put the right graph's Koszul class in a sum of tensor patterns
with at most two active factors outside the multiplication kernels.
A compatible projection for the induced exterior differential then gives
pure tensors of classes with at least \(m-2\geq3\) active factors both
closed for that differential and killed by \(b_i\).
These give three independent parameter lifts, and the source triple
membership forces their pairing to vanish.  The closure statements concern
Koszul classes; the chosen form representatives need not be closed for
ordinary \(\dd\).

\subsection{Algebraic approximation and finite specialization}

\begin{lemma}\label{lem:tensor-reduction}
If the data in Theorem~\ref{thm:tensor-obstruction} existed, data satisfying
the same conditions would exist over a field of positive characteristic
\(p\), for arbitrarily large \(p\).  The edge functions would still be
separated in their original blocks.
\end{lemma}

\begin{proof}
First put the parametrizations into graph form.  Choose complementary
linear projections adapted to their central tangent spaces.  For \(L\),
the chosen projection of \(H_L\) has invertible \(x\)-Jacobian modulo
\((x,t)\).  It therefore has a formal inverse over \(T\), with the parameters
fixed.  Its constant displacement is in \((t)\), so this is a legitimate
substitution in the complete local ring \(T[[x]]\).
For \(R\) the same assertion is the ordinary formal inverse function
theorem.  After these reparametrizations each graph has a fixed linear
projection equal to its independent variables.

We describe explicitly the approximation problem for \(L\).
For \(A\subset I\), let \(t_A=\prod_{i\in A}t_i\), including \(t_\varnothing=1\).
Expand its graph functions and all the coefficients witnessing
\eqref{eq:tensor-source-membership} in this finite basis:
\[
 H_L(x,t)=\sum_{A\subset I}t_A H_A(x),\qquad
 c(x,t)=\sum_{A\subset I}t_A c_A(x).
\]
All unknown coefficient functions have the same list \(x\) of independent
variables.

The fixed algebraic functions \(F_i,b_i\) admit finite pointed étale
presentations.  In such a presentation, differentiation in \(h_\nu\)
extends uniquely from the polynomial coordinate ring: implicit
differentiation expresses the derivative of every represented function
using the inverse étale Jacobian.  Thus the functions and their ambient
first derivatives can be handled in one finite presentation.  Introduce
auxiliary unknowns for its coordinates and its inverse units, and expand
them too in the basis \(t_A\).  Equating coefficients in the following
identities gives a finite polynomial system:
\begin{itemize}
 \item the étale presentation equations and the inverse-unit equations;
 \item the fixed graph-projection equation;
 \item \(S_t(H_L)=0\);
 \item each equation expressing \(t_i t_j t_k\) as a linear combination
       of the restricted ambient gradient entries.
\end{itemize}
No derivative of an unknown graph function occurs.  The only derivatives
are the fixed ambient derivatives of the given potentials, represented
as just described.  The existing formal graph and membership witnesses
give a formal solution of this polynomial system.

Artin approximation over the henselization \(K[x]_{(x)}^h\)
\cite[Theorem~1.10]{Artin1969} supplies an algebraic solution preserving
the desired finite jets.  Preserve the centers, the branch centers of
the étale presentations, and the central first graph jet.
Do the same independently for \(R\), encoding
\eqref{eq:tensor-target-membership} in its own independent variables \(y\).
These two systems share only fixed functions.  Their compatibility
condition is the open condition of complementary central tangents, which
the preserved jets retain.  There is consequently no nested approximation
problem.  If one independent-variable list is empty, the corresponding
system is already a finite system over \(K\).

We now have finitely many algebraic functions and finitely many exact
identities in henselizations.  Each function is represented on a finite
pointed étale neighborhood.  An identity in a henselization holds at some
finite stage of the filtered system of these neighborhoods.  Enlarging
the stages finitely many times represents all the functions, compositions,
membership identities, branch centers, and inverse units together.
Their finite presentations descend to a finitely generated integral
\(\mathbb Z\)-subalgebra \(A\subset K\), after adjoining the required
inverse units.  In particular we retain invertibility of the étale
Jacobians and of the transverse tangent determinant.  The implicit
differentiation formulas descend as well, so the specialized gradient
entries remain the derivatives of the specialized potentials.
The original presentations for the edge functions involve only their
own blocks \(h_i\), and retain that property after descent.

For completeness, this finite localization still has closed points in
arbitrarily large positive characteristics.  The nonzero finite-type
\(\mathbb Q\)-algebra \(A\otimes_{\mathbb Z}\mathbb Q\) has a maximal
ideal with residue field a number field \(E\).  The images of the finitely
many generators of \(A\), including all inverted elements, belong to
\(\mathcal O_E[1/d]\) for some positive integer \(d\).
Reduction at any prime of \(\mathcal O_E\) over a rational prime not
dividing \(d\) gives a homomorphism from \(A\) to a finite field.
All the declared units remain units.  The reduced pointed étale
presentations give formal germs over that residue field, satisfying the
same identities and the same tangent condition.
\end{proof}

\begin{remark}\label{rem:tensor-coefficient-field}
The lemma also applies when \(K=\mathbb C((\sigma))\).
Only finitely many elements of \(K\) occur as coefficients of the finite
presentations, and each is treated as a single coefficient.
Specialization is performed on \(A\), not coefficientwise on Laurent
series or on the whole field \(K\).
In particular, no boundedness of the negative \(\sigma\)-orders of the
coefficients of an Artin solution is required.
The algebraicity of the fixed separated potentials is essential:
the argument does not assert that ordinary Artin approximation can
algebraize arbitrary unknown formal potentials while preserving separate
variable restrictions.
\end{remark}

\subsection{Truncated forms and transverse graph classes}

The approximation step preserves all graph identities and actual ideal
memberships. We now turn their positive-characteristic specializations
into identities in finite-dimensional complexes, where transverse
graphs can be tested by a nonzero pairing.

Fix a specialization supplied by Lemma~\ref{lem:tensor-reduction} in
characteristic \(p>2\), and denote its field by \(\kk\).
Continue to use the same letters for the specialized germs.
Write \(z_1,\ldots,z_N\) for all internal coordinates and set
\[
 A_p=\kk[z_1,\ldots,z_N]/(z_1^p,\ldots,z_N^p),\qquad
 \Omega_p=A_p\otimes_\kk\bigwedge\nolimits_\kk
                   \langle \dd z_1,\ldots,\dd z_N\rangle .
\]
This is the ordinary algebra of differential forms of \(A_p\).
Its differential is well defined because \(\dd(z_\nu^p)=0\).
With external parameters, use
\[
 T_p=\kk[t_i:i\in I]/(t_i^2:i\in I),
 \qquad \Omega_{p,T}=\Omega_p\otimes_\kk T_p,
\]
and differentiate only in the internal coordinates.
All these spaces are finite over the indicated coefficient rings.

\begin{lemma}\label{lem:tensor-frobenius}
Every formal change of internal coordinates over \(T_p\) that is centered
modulo \((t)\) preserves the ideal
\((z_1^p,\ldots,z_N^p)\).  Hence it induces an automorphism of the truncated
algebra and of its relative differential forms.
\end{lemma}

\begin{proof}
If \(c\in(z,t)T_p[[z]]\), Frobenius shows that \(c^p\) belongs to
\((z_1^p,\ldots,z_N^p)\): every nonconstant internal monomial acquires an
exponent divisible by \(p\), and every positive parameter monomial has
zero \(p\)-th power.  Apply this to each new coordinate.
The inverse coordinate change has the same property, giving equality of
the ideals and the asserted automorphisms.
\end{proof}

Define a \(\kk\)-linear functional on forms, zero outside exterior degree
\(N\), by
\[
 \int f(z)\,\dd z_1\wedge\cdots\wedge\dd z_N
   =[z_1^{p-1}\cdots z_N^{p-1}]f .
\]
Extend it \(T_p\)-linearly in the presence of parameters.
It vanishes on exact forms: a derivative contributing exponent \(p-1\)
in its differentiated variable would have to come from exponent \(p\),
whose derivative is zero.  Consequently, for homogeneous forms,
\begin{equation}\label{eq:tensor-integration-by-parts}
 \int\dd\alpha\wedge\beta
       =-(-1)^{|\alpha|}\int\alpha\wedge\dd\beta .
\end{equation}

Let \(c_1,\ldots,c_r\) be graph equations for a smooth formal graph over
the parameter ring.  Their joint center is zero modulo \((t)\).
Its graph form is
\begin{equation}\label{eq:tensor-thom-form}
 U_c=\bigl(c_1^{p-1}\dd c_1\bigr)\wedge\cdots\wedge
                       \bigl(c_r^{p-1}\dd c_r\bigr)
       \in\Omega_{p,T}^{r}.
\end{equation}
This expression involves only a finite jet after truncation.

\begin{lemma}\label{lem:tensor-graph-form}
The form \(U_c\) is \(\dd\)-closed and is annihilated by the graph ideal.
Every internal one-form whose restriction to the graph is zero wedges
to zero with \(U_c\).
For two central graphs with complementary tangents and graph forms
\(U_0,V\),
\begin{equation}\label{eq:tensor-nonzero-pairing}
                     \int U_0\wedge V\ne0 .
\end{equation}
\end{lemma}

\begin{proof}
Each factor \(c_j^{p-1}\dd c_j\) is closed, and \(c_j^p=0\) in the
truncated ring by Lemma~\ref{lem:tensor-frobenius}.
This proves the first two assertions.
In graph coordinates, the kernel of restriction of internal one-forms is
the sum of graph-ideal multiples of one-forms and the span of the
\(\dd c_j\).  Both summands wedge to zero with \(U_c\).
This description remains valid after truncation because the graph
coordinate change preserves the truncation ideal.

For the pairing, concatenate the central graph equations for the two
graphs, and let \(M\) be their linear coefficient matrix in the original
coordinates.  Complementarity says that \(M\) is invertible.
The coefficient part of \(U_0\wedge V\) has degree at least \(N(p-1)\);
all terms of higher degree vanish in \(A_p\).
Thus only the linear terms of the graph equations and the constant terms
of their differentials contribute.  A linear coordinate change \(M\)
acts on the polynomial socle
\(\kk z_1^{p-1}\cdots z_N^{p-1}\) by \(\det(M)^{p-1}\), and on the
exterior top form by \(\det(M)\).  To check the socle assertion one may
use diagonal matrices, permutations, and elementary transvections:
for a transvection every nonconstant binomial term has an exponent at
least \(p\) and vanishes.  These matrices generate the general linear
group.  The pairing is therefore \(\det(M)^p\), up to the choice of order
of the graph equations, and is nonzero.
\end{proof}

Let \(U_t\) be the graph form for \(L\), and let \(V\) be the graph form
for \(R\).  Define the degree-one operators
\[
       D_t=\dd S_t\wedge(-),\qquad D=\dd S_0\wedge(-).
\]
Both square to zero.  The graph identities and
Lemma~\ref{lem:tensor-graph-form} give
\begin{equation}\label{eq:tensor-graph-cycles}
 D_tU_t=0,\qquad DV=0,\qquad
 \dd U_t=0,\qquad\dd V=0 .
\end{equation}
In particular \(DU_0=0\) and the nonzero pairing
\eqref{eq:tensor-nonzero-pairing} is defined.

\begin{lemma}\label{lem:tensor-membership-boundary}
For every three distinct active indices \(i,j,k\), there are actual forms
\(B_{ijk}\) and \(C_{ijk}\) with
\begin{equation}\label{eq:tensor-actual-boundaries}
 t_i t_j t_k U_t=D_t B_{ijk},\qquad
 b_i b_j b_k V=D C_{ijk}.
\end{equation}
The first equality is over the full parameter algebra \(T_p\).
\end{lemma}

\begin{proof}
Lift the graph membership coefficients to the ambient ring.
The difference between the asserted scalar and the resulting ambient
gradient combination belongs to the graph ideal, so it annihilates the
graph form.  If \(\iota_\nu\) is contraction with the coordinate vector
\(\partial/\partial z_\nu\), then
\[
 \iota_\nu D_t+D_t\iota_\nu
                =(\partial_{z_\nu}S_t)\,\id.
\]
For a \(D_t\)-cycle \(U_t\), multiplication by
\(\sum_\nu a_\nu\partial_{z_\nu}S_t\) is consequently
\[
 \left(\sum_\nu a_\nu\partial_{z_\nu}S_t\right)U_t
       =D_t\left(\sum_\nu a_\nu\iota_\nu U_t\right).
\]
The operator \(D_t\) does not differentiate the coefficients \(a_\nu\).
This proves the first equality, and the same argument with \(D,V\)
proves the second.
\end{proof}

We will use two elementary facts about this pairing.  If
\(\theta\) is a one-form and \(D_\theta=\theta\wedge(-)\), then, for
homogeneous \(\alpha\),
\begin{equation}\label{eq:tensor-wedge-adjoint}
 (D_\theta\alpha)\wedge\beta
          =(-1)^{|\alpha|}\alpha\wedge D_\theta\beta.
\end{equation}
Thus a \(D_\theta\)-boundary wedges to zero with a
\(D_\theta\)-cycle, even before integration.
Also \(\dd D+D\dd=0\), since \(\dd^2S_0=0\).
Ordinary differentiation therefore induces a differential on the
homology of \(D\); Equation~\eqref{eq:tensor-integration-by-parts}
continues to test its boundaries against actual closed representatives.

\subsection{Two homologies and a compatible projection}

We have graph cycles with a nonzero pairing and actual boundary
identities for every triple. The separate edge blocks now let us
analyze the target graph class one factor at a time. We first take
Koszul homology, then use the differential induced on it by ordinary
exterior differentiation.

Let \(\Omega_i\) be the truncated form algebra in the block \(h_i\), and
write
\[
 D_i=\dd F_i\wedge(-),\qquad H_i=H(\Omega_i,D_i).
\]
The tensor products of graded spaces and operators in what follows use
the usual Koszul signs.  The block separation gives
\[
 (\Omega_p,D)=\bigotimes_{i=1}^n(\Omega_i,D_i),\qquad
 H(\Omega_p,D)=\bigotimes_{i=1}^n H_i .
\]
The second equality is the finite-dimensional Künneth isomorphism over
the field \(\kk\).
Let \(\delta_i\) denote the differential induced on \(H_i\) by ordinary
\(\dd\), and let \(\delta\) be their tensor differential.
For \(i\in I\), multiplication by \(b_i\) induces a degree-zero
endomorphism \(B_i\) of \(H_i\), since it commutes with \(D_i\).
Set
\[
                         K_i=\ker B_i .
\]
There is no assertion that \(K_i\) is a \(\delta_i\)-subcomplex.

\begin{lemma}[Kernel support]\label{lem:tensor-kernel-support}
Let \(v\in\bigotimes_{i=1}^nH_i\) be killed by
\(B_iB_jB_k\) for every three distinct active indices.
Then \(v\) belongs to the sum of tensor subspaces having all but at most
two active factors in the corresponding \(K_i\).
\end{lemma}

\begin{proof}
Choose a graded complement \(J_i\) to \(K_i\) in each active factor.
The restriction of \(B_i\) to \(J_i\) is injective.
For a fixed triple, the kernel of its tensor operator is the sum of
subspaces having a \(K_i\) factor in at least one of the three specified
positions: the induced tensor of injections on the quotient by those
kernels is injective.
Expand the whole tensor product using \(H_i=K_i\oplus J_i\).
Intersecting the kernel conditions for all triples leaves exactly the
summands with at most two active \(J_i\) positions.
\end{proof}

\begin{lemma}[Projection preserving kernel representatives]
\label{lem:tensor-kernel-projection}
Let \((H,\delta)\) be a finite-dimensional graded complex over a field,
and let \(K\subset H\) be any graded subspace.
There is a chain projection \(P:H\to H\), homotopic to the identity,
whose image is a space of representatives for \(H(H,\delta)\) and which
satisfies
\[
                    P(K)\subset K\cap\ker\delta .
\]
\end{lemma}

\begin{proof}
Write \(Z=\ker\delta\) and \(B=\im\delta\).
Choose a complement \(A_K\) to \(K\cap B\) in \(K\cap Z\), and extend
\(A_K\) to a complement \(A\) to \(B\) in \(Z\).
Choose a complement \(W_K\) to \(K\cap Z\) in \(K\).
It is disjoint from \(Z\), so extend it to a complement \(W\) to \(Z\)
in \(H\).  All these choices are made degree by degree.
Then
\[
                       H=B\oplus A\oplus W .
\]
Let \(P\) be projection onto \(A\).  It is a chain map, and
\(P(K)\subset A_K\subset K\cap Z\).
The map \(\delta:W\to B\) is an isomorphism with the appropriate shift
of degrees.  Define a degree-minus-one map \(s\) to be its inverse on
\(B\), and zero on \(A\oplus W\).  Then
\[
                        \id-P=\delta s+s\delta .
\]
\end{proof}

Apply Lemma~\ref{lem:tensor-kernel-projection} to each active factor
\((H_i,\delta_i,K_i)\), and choose an ordinary homology projection in
each inactive factor.  Denote these maps by \(P_i\), and put
\(\Pi=\bigotimes_iP_i\).  This tensor projection is chain-homotopic
to the identity.  Indeed, the telescoping identity for
\(\id-\bigotimes_iP_i\), together with the individual homotopies and
the graded tensor convention, supplies a homotopy for the total
differential \(\delta\).

Let \(v=[V]\in H(\Omega_p,D)\).  Equations
\eqref{eq:tensor-actual-boundaries} imply the hypotheses of
Lemma~\ref{lem:tensor-kernel-support}, and \(\delta v=0\) by
\eqref{eq:tensor-graph-cycles}.  Thus \(\Pi v\) is a sum of homogeneous
pure tensors
\begin{equation}\label{eq:tensor-pure-cycles}
                   v_1\otimes\cdots\otimes v_n
\end{equation}
such that every \(v_i\) is a \(\delta_i\)-cycle and at least \(m-2\)
active factors lie in \(K_i\cap\ker\delta_i\).
This conclusion follows by applying \(\Pi\) to the tensor subspaces in
Lemma~\ref{lem:tensor-kernel-support}; the projection of a kernel factor
is still a kernel factor.
Since \(m\geq5\), every nonzero pure tensor in this sum has at least three
such active factors.

Choose \(D_i\)-cycle representatives for all the factors in
\eqref{eq:tensor-pure-cycles}, and let \(V'\) be the resulting sum of
actual forms.  The chain homotopy for \(\Pi\) shows that
\([V']-[V]=\delta[Z]\) for some class in \(H(\Omega_p,D)\).
Choose a \(D\)-cycle representative \(Z\).
There is then an actual form \(E\) such that
\begin{equation}\label{eq:tensor-chain-replacement}
                         V'-V=\dd Z+D E .
\end{equation}
Pairing with \(U_0\), equations
\eqref{eq:tensor-integration-by-parts},
\eqref{eq:tensor-graph-cycles}, and
\eqref{eq:tensor-wedge-adjoint} give
\begin{equation}\label{eq:tensor-pairing-preserved}
                    \int U_0\wedge V'=\int U_0\wedge V .
\end{equation}
The representatives composing \(V'\) need not be closed for ordinary
\(\dd\).  Their classes are \(\delta_i\)-cycles, which is precisely what
was used to obtain \eqref{eq:tensor-chain-replacement}.

\subsection{Lifts over a three-parameter cube}

The replacement class has the same pairing with the source graph as
the original target class. Each of its tensor summands has at least
three active factors lying in both the multiplication kernel and the
kernel of the induced differential. Those three factors supply the
parameter lifts used to force that summand's pairing to vanish.

\begin{proof}[Proof of Theorem~\ref{thm:tensor-obstruction}]
Suppose the stated data exist and perform the reduction and constructions
above.  Consider one pure tensor \eqref{eq:tensor-pure-cycles} contributing
to \(V'\).  Choose three active positions \(i,j,k\) in which
\[
                   B_\nu v_\nu=0,\qquad\delta_\nu v_\nu=0 .
\]
On \(H_\nu\), the commutator identity is
\[
       [\delta_\nu,B_\nu]
         =\delta_\nu B_\nu-B_\nu\delta_\nu
         =[\,\dd b_\nu\wedge(-)\,].
\]
It follows that \([\,\dd b_\nu\wedge(-)\,]v_\nu=0\).
For a homogeneous \(D_\nu\)-cycle representative \(a_\nu\) of \(v_\nu\),
there is therefore a form \(w_\nu\) of the same exterior degree with
\[
                    D_\nu w_\nu=-\dd b_\nu\wedge a_\nu .
\]
Consequently
\begin{equation}\label{eq:tensor-first-order-lift}
 (D_\nu+t_\nu\dd b_\nu\wedge)
                      (a_\nu+t_\nu w_\nu)=0
             \quad\text{in }\Omega_\nu[t_\nu]/(t_\nu^2).
\end{equation}
Only the vanishing of the indicated classes in \(H_\nu\) was required;
no ordinary \(\dd\)-closed representative is being assumed.

Set all parameters outside this triple equal to zero, and put
\[
 T_{ijk}=\kk[t_i,t_j,t_k]/(t_i^2,t_j^2,t_k^2).
\]
Tensor the three lifts \eqref{eq:tensor-first-order-lift} with the chosen
unchanged \(D_\nu\)-cycle representatives in every other block.
The resulting form \(W_t\) is a cycle for
\[
              \dd(S_0+t_i b_i+t_j b_j+t_k b_k)\wedge(-).
\]
This is an exact cycle equation over \(T_{ijk}\): each perturbation acts
in its own block, so the three separate first-order lifts require no
additional mixed corrections.

Specialize the actual first identity in
\eqref{eq:tensor-actual-boundaries} to these three parameters.
It remains an actual boundary identity in forms.  In particular, this
step does not assert that homology commutes with the possibly nonflat
parameter specialization.
By \eqref{eq:tensor-wedge-adjoint}, pairing that identity with \(W_t\)
gives
\[
 t_i t_j t_k\,P(t)=0\quad\text{in }T_{ijk},\qquad
 P(t)=\int U_t\wedge W_t ,
\]
where \(U_t\) is specialized in the same way.
The square-free monomials form a \(\kk\)-basis of \(T_{ijk}\).
Multiplication by \(t_i t_j t_k\) kills all its positive-degree
parameter monomials, so the displayed equation implies
\[
  P(0)=\int U_0\wedge(a_1\otimes\cdots\otimes a_n)=0 .
\]
The chosen triple may vary with the pure tensor; the argument applies
to each term separately.  Summing gives \(\int U_0\wedge V'=0\).
This contradicts \eqref{eq:tensor-pairing-preserved} and the nonzero
transverse graph pairing \eqref{eq:tensor-nonzero-pairing}.
\end{proof}
